\documentclass[../main.tex]{subfiles}
\begin{document}
%==========================================TIÊU ĐỀ TẬP========================================
\begin{table}[h]
  \begin{adjustwidth}{-1cm}{}
    \begin{tabular}{>{\centering\arraybackslash}p{6cm}|>{\raggedright\arraybackslash}p{12.5cm}}
      \multirow{5}{6cm}
      % Cột bên trái: ÉPISODE và số tập
      {\\[-40pt]
        \flaregothic\fontsize{20pt}{20pt}\selectfont \centering
        {\contour{errorfive}{\textcolor{black}{ÉPISODE}}}\color{black}\\
        \burbank\fontsize{72pt}{72pt}\selectfont
        \includegraphics[height=3.12359cm]{../../../../Image/Book?/number/num34_alt.png}
        \\[18pt]
        % \flaregothic\fontsize{14pt}{14pt}\selectfont
        % \color{epattr}BIENVENUE, \uline{\color{Banpire} Mel} !
        \flaregothic\fontsize{14pt}{14pt}\selectfont{\contour{errorfive}{\textcolor{white}{BIENVENUE, }}
          {\color{black} Naomi}
          {\contour{errorfive}{\textcolor{white} {!}}}}\\
        % \sen\fontsize{18pt}{18pt}\selectfont{\contour{errorfive}{\textcolor{white}{Mini-histoire}}}
        % \color{black}
      }
       &
      % Dòng thứ nhất: tên tập tiếng Nhật
      {\fotskip\fontsize{18pt}{18pt}\selectfont \ruby{嵐}{あらし}の\ruby{後}{あと}の\ruby{静}{しず}けさ}
      \\[6pt]
      % Dòng thứ hai: tên tập tiếng Anh
       & {\cafeteria\fontsize{18pt}{18pt}\selectfont The Calm After the Storm}                                     \\[4pt]
      % Dòng thứ ba: tên tập tiếng Việt
       & {\vnmsans\fontsize{18pt}{18pt}\selectfont Sự bình yên sau cơn bão}                                        \\[8pt]
      % Dòng thứ tư: nhân vật xuất hiện
       & {\caviar{\fontsize{14pt}{14pt}\selectfont Nhân vật: \emp{kuon1}, \emp{kaigou1}, \emp{suzuki1}, \emp{naomi1}}}
      % \\[6pt]
      % Dòng cuối cùng: Thời gian diễn ra của tập (thường sẽ là ngày phát hành)
      % & {\continuum\fontsize{16pt}{16pt}\selectfont Ngày phát hành: 11/06/2022}\\[8pt]
    \end{tabular}
  \end{adjustwidth}
\end{table}
%===========================================NỘI DUNG==========================================
\color{black}
\vspace{30pt}

\pov{Hikawa Kuon}

``\vie{\dots-chan.}''

``\vie{\dots Hửm?}''

``\vie{\dots Onii-chan, dậy đi nào.}''

Một giọng nói vang lên đâu đó.

``\vie{\dots Cái gì vậy\dots?}''

Hôm nay là ngày nghỉ cuối tuần. Theo thông lệ, vì không phải đến trường, cả ba chúng tôi thường cho phép mình nướng thêm một chút. Tôi thường chỉ ra khỏi giường vào khoảng bảy, tám giờ sáng. Thế nhưng, hôm nay lại khác.

Tôi lờ mờ tỉnh giấc bởi một cảm giác kỳ lạ, thứ mà một kẻ đã quen với việc ngủ một mình như tôi ít khi trải nghiệm. Có một sức nặng nào đó đang đè lên ngực tôi, không đến mức ngộp thở, nhưng đủ để cảm nhận được sự hiện diện của một thực thể vật lý.

Hơn nữa, vật thể ấy còn tỏa ra một hơi ấm dìu dịu. Đó chắc chắn không thể là Kaigou hay Suzuki, bởi vì với vóc dáng của hai cô nàng đó, tôi sẽ bị đè bẹp dí chứ không phải là cái sức nặng nhẹ bẫng như thế này.

Và rồi, giọng nói kia lại vang lên ngay bên tai tôi. Lần này lại sát hơn.

``\vie{Onii-chan, trời sáng rồi đó.}''

Tôi hơi khẽ nhíu mày khi mắt vẫn nhắm nghiền. ``Onii-chan'' ư? Cách gọi này\dots\ làm tôi nhớ đến Ryuuhou, đứa em gái nhỏ của tôi ở thế giới thực. Nhưng đó là cách con bé gọi tôi hồi còn nhỏ, trước khi bước vào tuổi dậy thì và chuyển sang kiểu gọi ``Onii-san'' đầy xa cách (tôi cũng chẳng để tâm lắm tới chuyện đó; dù gì tôi cũng có phải là một tên cuồng em gái đâu). Hơn nữa, giọng nói này không phải của Ryuuhou. Nó trầm hơn, nhẹ nhàng hơn, và mang một sắc thái tĩnh lặng lạ thường, như tiếng lá rơi trong một khu rừng vắng.

Tôi chầm chậm mở mắt.

``\dots''

``\vie{Ồ, anh dậy rồi kìa.}''

Đập vào mắt tôi, ngay trong cự ly gần nhất có thể, là một khuôn mặt nhỏ nhắn. Một mái tóc đen được tết gọn gàng thành hai bím nhỏ vắt qua vai, làm nổi bật lên nước da trắng ngần. Nhưng điều khiến tôi chú ý nhất là đôi mắt. Đôi mắt màu xanh lục bảo trong vắt, sâu thẳm như chứa đựng cả một nghìn năm sương khói của rừng Aogami. Cô bé đang khoác trên mình một bộ kimono mộc mạc, tĩnh lặng nằm úp trên ngực tôi, chống cằm nhìn tôi với một sự hiếu kỳ không che giấu, nhưng lại vô cùng trầm mặc.

Khoan đã.

Tôi chớp mắt vài cái. Bộ não ngái ngủ của tôi bắt đầu khởi động. Mắt xanh lục\dots\ tóc bím\dots\ kimono\dots

Đây chẳng phải là oán linh ``hạt nhân'' của rừng Aogami sao!? Tại sao lại ở trong phòng ngủ của tôi?

Tôi liếc mắt nhìn chiếc đồng hồ trên tủ đầu giường. Sáu giờ sáng. Giờ này là giờ âm dương giao thoa hay gì? Mình lại đang nằm mơ thấy ma à?

Nhưng khi cô bé khẽ chớp mắt, và một lọn tóc của em chạm vào má tôi, mang theo cảm giác nhồn nhột rất thực tế\dots\ tôi mới nhận ra đây hoàn toàn không phải là một giấc mơ.

``{\Large\textbf{Ối giời ơi!}}''

Tôi bật ngửa - hay đúng hơn là cố bật ngửa, nhưng vì bị đè nên chỉ tạo ra một cú vặn mình vụng về, vô tình hất cô bé ngã lăn sang một bên nệm. Tiếng hét thất thanh của tôi phá tan bầu không khí yên tĩnh của buổi sáng, đủ lớn để đánh thức hai cái xác ướp đang cuộn tròn trong chăn ở giường bên kia.

``\vie{Cái gì!? Động đất à!?}'' Kaigou lồm cồm bò dậy, tóc tai bù xù như một tổ chim, tay vẫn còn nắm chặt cái gối.

``\vie{Ưm\dots\ Onii-sama\dots\ có chuyện gì\dots}'' Suzuki dụi mắt, ngáp dài một cái, giọng vẫn còn ngái ngủ.

Nhưng rồi, cả hai ngay lập tức hóa đá khi ánh mắt họ chạm phải sinh vật đang ngồi chễm chệ trên giường của tôi. Cô bé mắt xanh xoa xoa chỗ vừa bị ngã, đôi mắt lục bảo chớp chớp nhìn Kaigou và Suzuki bằng vẻ tĩnh lặng cố hữu.

``\vie{\textbf{M-Ma!!}}'' Kaigou hét lên, túm lấy chăn kéo lên tận cổ.

``\vie{Không, khoan đã Onee-sama,}'' Suzuki nhanh chóng nhận ra điểm bất thường, đôi mắt mở to. ``\vie{Ma không có bóng, cũng không chạm vào nệm lõm xuống như vậy được! Cô bé\dots\ cô bé là người sống!}''

Sự thật đó giáng xuống đầu ba chúng tôi như một gáo nước lạnh. Hàng loạt câu hỏi thi nhau bùng nổ trong đầu tôi. Tại sao cô bé lại ở đây? Tại sao một oán linh bị giam cầm 1300 năm lại trở thành một con người bằng xương bằng thịt? Và quan trọng nhất: làm thế nào mà em ấy vào được đây, lại còn nói được ngôn ngữ của chúng tôi?

``\vie{Mọi người trông có vẻ hoảng ghê\~{}}'' cô bé mắt xanh nói ngân dài, biểu cảm ngây thơ nhìn chúng tôi.

``\vie{D-Dĩ nhiên là hoảng rồi!}'' Tôi chỉ tay vào cô bé, giọng vẫn còn chút hoảng hốt. ``\vie{Làm thế nào mà nhóc lại vào được đây?}''

Cô bé không nói gì, chỉ chậm rãi đưa tay chỉ về phía cửa sổ ban công. Nhìn theo hướng đó, tôi mới nhận ra rằng cửa sổ ban công vẫn đang mở hé đón gió, chỉ có tấm rèm là được kéo lại. Có lẽ đêm qua tôi đã quên chốt khóa cửa ban công lại chăng?

``\vie{Được rồi, tạm gác chuyện an ninh nhà cửa lại,}'' Kaigou vuốt mặt, cố gắng bình tĩnh. ``\eng{Giải thích cho tôi hiểu, tại sao con nhóc oán linh lại đang ngồi đây, thở ra khí cacbonic và chiếm chỗ trên giường của chúng ta? Game-san, ra mặt giải thích đi!}''

Tôi chạm vào mặt đồng hồ. Một tia sáng xanh nhạt lóe lên, và giọng nói cơ khí đều đều của Game vang lên trong phòng: ``\eng{Dựa trên những gì tôi thu thập được từ tàn dư của nghi lễ đêm đó, đây là kết quả của cơ chế `Bảo tồn Dữ liệu'. Khi oán khí của rừng Aogami bị thanh tẩy hoàn toàn, linh hồn `hạt nhân' - tức cô bé này - vốn dĩ phải tan biến theo. Tuy nhiên, nhờ vào việc thiết lập một sự cộng hưởng vật lý và tâm linh cực mạnh với người thực hiện vai trò `Kết nối' - tức Suzuki-user - trước khi sự phân rã xảy ra, hệ thống thế giới đã vô tình ghi nhận dữ liệu linh hồn của thực thể này như một phần của thực tại vật lý mới.}''

Tôi khẽ nuốt nước bọt, cảm thấy một luồng điện lạnh chạy dọc sống lưng. Nếu cảm xúc của Suzuki đủ mạnh để buộc thế giới phải ``lưu lại'' một linh hồn, thì chúng tôi đang đùa với thứ gì đó lớn hơn trí tưởng tượng của mình rất nhiều.

``\eng{Cho cậu đang thắc mắc thì, tôi cũng đã dùng cơ chế đó để tái tạo lại Aogami và hồi sinh tất cả mọi người ở đó đấy,}'' Game tiếp tục.

``\eng{Ông có thể làm thế ư?}'' tôi nhíu mày hỏi.

``\eng{Duh, tôi là một chương trình mà. Bây giờ cậu mới biết khả năng của tôi à?}''

``\eng{Thế còn việc ngôn ngữ thì sao?}'' tôi vặn vẹo tiếp, liếc sang cô bé đang tĩnh lặng nhìn chúng tôi. ``\eng{Nhóc này là người từ cách đây một nghìn ba trăm năm cơ mà. Sao lại nói ngôn ngữ của chúng tôi thành thạo như người bản xứ vậy?}''

``\eng{Đó cũng là một phần của sự `Cộng hưởng'. Quá trình liên kết tâm trí với Suzuki-user không chỉ truyền tải cảm xúc, mà còn trực tiếp sao chép toàn bộ hệ thống ngôn ngữ và khái niệm thế giới quan hiện đại vào bộ nhớ của cô bé. Cứ tưởng tượng như cô bé vừa được tải về và cài đặt một gói ngôn ngữ tự động,}'' Game từ tốn đáp. ``\eng{Cô bé vẫn có thể nói được tiếng Nhật, nên cậu chớ lo.}''

Nói một cách dễ hiểu hơn: Suzuki đã trở thành chiếc mỏ neo kéo linh hồn của cô bé lại thế giới này, đồng thời là ``từ điển sống'' biến em thành một con người thực sự của thời đại này thay vì để em siêu thoát. Giống như cách mà tôi trở thành chỗ dựa vững chắc mà hai chị em họ có thể dựa dẫm sau khi từ bỏ thế giới cũ. Suzuki có lẽ là người hiểu rõ hơn hết cái cảm giác đó.

``\vie{Là\dots\ do em muốn ở lại sao?}'' cô em gái mắt hai màu tiến lại gần, dịu dàng hỏi cô bé.

Cô bé mắt xanh khẽ gật đầu, giọng nói trầm tĩnh cất lên: ``\vie{Em\dots\ không muốn bị bóng tối nuốt chửng nữa. Khi bàn tay của chị chạm vào em, em cảm nhận được sự ấm áp. Em muốn nó. Thế nên, thế giới đã nghe thấy và cho em một thân xác\dots\ để em không bị tan biến.}''

Lời giải thích mang đậm màu sắc tâm linh của cô bé hòa quyện một cách hoàn hảo với sự khô khan kỹ thuật của Game. Dù hiểu theo cách nào, thì sự thật trước mắt vẫn không thể thay đổi: chúng tôi vừa có thêm một nhân khẩu trong căn hộ nhỏ này.

\ct

Sau một hồi dỗ dành và sắp xếp lại những thông tin, chúng tôi ngồi quây quần quanh bàn tròn ở phòng khách, trong khi cô bé đang tò mò nhấm nháp ly sữa ấm mà Suzuki vừa pha.

``\vie{Vậy\dots\ em có tên không?}'' tôi hỏi, chống cằm nhìn cô bé. ``\vie{Game có nói em từng có một `định danh' ở thời đại cũ đúng không?}''

Cô bé đặt ly sữa xuống, đôi mắt lục bảo nhìn thẳng vào tôi. ``\vie{Cái tên đó thuộc về người đã chết từ một nghìn ba trăm năm trước rồi ạ. Em không muốn dùng lại nó. Em muốn\dots\ sống một cuộc đời mới. Xin các anh chị\dots\ hãy đặt cho em một cái tên.}''

Một cái tên mới cho một khởi đầu mới. Điều này hoàn toàn hợp lý. Touka cũng đã từng nói vậy khi giải thích về ngày Lễ tình nhân rơi vào tháng Năm thay vì tháng Hai như ta biết.

``\vie{Vậy\dots\ Midori thì sao? Mắt em màu xanh lục mà?}'' Kaigou đề xuất ngay lập tức, một cái tên kinh điển không thể đơn giản hơn.

Cô bé hơi nghiêng đầu, đôi mắt lộ vẻ không cảm xúc như kiểu: ``Chị đang đùa em đấy à?'' Kaigou lập tức ho khan một tiếng, rút lại ý kiến. ``\vie{B-Bỏ đi. Coi như chị chưa nói gì.}''

``\vie{Thế còn\dots\ Shizuka? Vì em ấy rất điềm tĩnh và yên lặng,}'' Suzuki đưa ra một lựa chọn thanh nhã hơn.

Cô bé suy nghĩ một chút, rồi khẽ lắc đầu, tỏ ý không thực sự ưng bụng. Tên đó cũng hơi đại trà thật. Cả hai chị em quay về phía tôi, đồng thanh hỏi: ``\vie{Kuon-kun/Onii-sama có ý tưởng gì không (ạ)?}''

Tôi gõ gõ ngón tay lên bàn, nhìn vẻ mặt trầm mặc nhưng lại toát lên một sức sống âm thầm của cô bé. Một cái tên vừa đẹp, vừa mang ý nghĩa của sự chân thật và vẻ đẹp tự nhiên.

``\vie{\dots Naomi,}'' tôi cất lời. ``\vie{Naomi (\ruby{直}{なお}\ruby{美}{み}, \ruby{Na-ô}{Trực}-\ruby{mi}{Mỹ}) - Vẻ đẹp của sự chân thật. Em nghĩ sao?}''

Cô bé mắt xanh tròn xoe mắt nhìn tôi. Em không nói gì, chỉ lẳng lặng lẩm nhẩm cái tên ``Naomi'' vài lần trong miệng. Rồi, một nụ cười hiếm hoi nở trên môi em - một nụ cười mỉm nhỏ bé nhưng chân thực, xua tan đi vẻ liêu trai tĩnh lặng.

``\vie{Naomi. Em thích nó.}''

Tôi thở phào nhẹ nhõm, cảm thấy một tảng đá vừa lăn khỏi lồng ngực. Cái tên ``Naomi'' như đã được sinh ra để dành cho cô bé này vậy.

``\vie{Phù\dots\ đặt được tên rồi,}'' Kaigou xoa vai, thở dài. ``\vie{Cậu có khiếu đặt tên đấy, Kuon-kun. Tên Midori của tôi nghe như tên một con cá, còn Shizuka của Suzu thì lại quá phổ biến. Naomi của cậu thì vừa thanh nhã, vừa có chiều sâu.}''

Tôi gãi gãi má, có chút ngượng. ``\vie{À, tại tôi thấy em ấy\dots\ rất chân thật. Không giả tạo, không che giấu.}''

Suzuki nắm tay tôi, mắt long lanh. ``\vie{Onii-sama, em biết mà! Anh đúng có khiếu đặt tên đấy chứ!}''

Tôi lắc đầu, không nhịn được cười. Cảm giác căng thẳng từ nãy giờ như tan biến hết. Cô bé từ 1300 năm trước giờ đã có một cái tên mới. Một khởi đầu mới. Và thế là, vào một buổi sáng cuối tuần bình yên sau bao sóng gió, cô bé mắt xanh của rừng Aogami đã chính thức mang tên Naomi.

\ct

Ngay sau khi đã quen với việc có một cơ thể vật lý, tính cách của Naomi bắt đầu bộc lộ rõ hơn qua những hành động nhỏ. Dù Game có nói rằng cô bé đã được tiếp thu toàn bộ khái niệm về thế giới quan hiện đại từ Suzuki, nhưng điều đó dường như chỉ giống như việc đọc một cuốn bách khoa toàn thư mà chưa từng một lần trải nghiệm thực tế. Kinh nghiệm của em về thế giới này hoàn toàn bằng không. Chưa kể đến việc quan niệm về cuộc sống của một cô bé mười lăm tuổi chưa thể sâu sắc bằng một người lớn được. Nhưng với một cô bé vừa được tái sinh, chừng đó cũng quá đủ để thích nghi với lối sống thế kỷ XXI rồi.

Naomi bước những bước nhỏ, rón rén như một chú mèo lại gần chiếc tivi màn hình phẳng, nhẹ nhàng vươn ngón tay chạm vào bề mặt kính đen bóng. ``\vie{Chiếc hộp này\dots\ sao lại phẳng và lạnh thế?}'' em thì thầm, mắt long lanh tò mò.

Khi Kaigou bật công tắc, màn hình đột ngột sáng lên khiến em hơi giật mình lùi lại một bước, nhưng không hề hoảng hốt. Đôi mắt lục bảo mở to, nhìn chằm chằm vào những hình ảnh chuyển động bên trong chiếc hộp nhựa với một sự say mê tĩnh lặng. ``\vie{Người\dots\ người ở trong đó kìa!?}'' Naomi nghiêng đầu, đưa tay vẽ nhẹ viền khung viền màn hình. ``\vie{Họ nhỏ xíu, nhưng lại động đậy được.}''

Kaigou bật cười, chống nạnh: ``\vie{Đó chỉ là hình ảnh thôi, không phải người thật đâu.}''

``\vie{Hình ảnh\dots\ ư?}'' Naomi gật gù chậm rãi, như thể đang nhai từ chữ một cách cẩn thận.

Sau đó, em lại quay sang nghịch công tắc đèn, bật tắt vài lần, và đứng lặng thinh nhìn chiếc quạt trần quay đều đều. ``\vie{Ánh sáng\dots\ không cần lửa,}'' em thì thào, ngón tay khẽ gõ nhịp theo từng vòng quay của quạt. ``\vie{Và gió\dots\ không cần lá. Thế giới này kỳ diệu quá.}''

Kaigou nhún vai, vừa cài lại kẹp tóc vừa nói: ``\vie{Ờm, còn nhiều thứ kỳ hơn thế nữa. Sau này em sẽ thấy hết thôi.}''

Sự hiếu kỳ của em thể hiện rõ rệt, nhưng không hề ồn ào hay phá phách như những đứa trẻ tám tuổi bình thường, mà lại toát lên một nét trầm mặc cổ xưa.

``\eng{Nhân tiện thì,}'' Game bất ngờ lên tiếng từ mặt đồng hồ, ``\eng{tôi đã quét sinh trắc học của thực thể này. Khung xương và cấu trúc tế bào khớp với độ tuổi sinh học là 8 tuổi. Dựa trên những phân mảnh ký ức tôi thu thập được, cô bé thực sự đến từ thời kỳ đầu Nara, loanh quanh thế kỷ thứ VIII.}''

``\eng{Thế kỷ VIII? Vậy là đúng như những gì mà Sakura-san nói sao?}'' tôi vuốt cằm. ``\eng{Vậy thì chúng ta nên làm gì với ngày tháng năm sinh của em ấy đây? Nếu định làm giấy tờ tùy thân hay hồ sơ nhập học, không thể bỏ trống được đâu.}''

Dù ở thời điểm này, cả ba chúng tôi đều chưa hề có khái niệm hay dự định đưa Naomi cùng trở về thế giới thực, nhưng việc tồn tại ở Aogami cũng cần một vỏ bọc pháp lý, chẳng may nếu chúng tôi bị kẹt lại lâu hơn dự tính.

``\eng{Dữ liệu lịch sử từ 1300 năm trước quá mơ hồ để xác định ngày sinh chính xác,}'' Game phân tích. Tôi không quá ngạc nhiên với kết quả này. Đến cả hóa thạch khủng long người ta cũng chỉ ước lượng chính xác tới thời kỳ nào chứ không thể cụ thể đến thế này. ``\eng{Chính vì thế, tôi đề xuất sử dụng ngày diễn ra nghi lễ hôm đó làm ngày sinh. Dù sao đó cũng là ngày cô bé được `tái sinh'. Còn về năm sinh, bối cảnh hiện tại là năm 2021. Trừ đi 8 tuổi sinh lý, năm sinh của cô bé sẽ là 2013.}''

``\vie{Ngày hôm qua\dots\ 2013,}'' tôi gật gù ghi nhận. ``\vie{Hợp lý đấy.}''

Tôi quay sang nhìn Naomi, người đang lật xem một cuốn tạp chí trên bàn dù có vẻ em thích thú với chất liệu giấy láng mịn hơn là nội dung bên trong.  ``\vie{Mà này, Naomi.}''

Cô bé ngẩng lên, nghiêng đầu chờ đợi.

``\vie{Sao em lại gọi anh là `Onii-chan'?}'' tôi hỏi điều mà tôi vẫn luôn thắc mắc từ lúc mở mắt. ``\vie{Một người đến từ 1300 năm trước như thế, theo lý thuyết thì phải xưng hô kiểu cách hơn chứ?}''

Nghe vậy, Naomi từ từ đặt cuốn tạp chí xuống, hai tay chắp lại ở vạt áo kimono.

``\vie{Đáng lẽ\dots\ em định gọi anh là `Onii-sama'.}'' Cô bé nhẹ nhàng nói, chất giọng tĩnh mịch đặc trưng. ``\vie{Nhưng chị ấy,}'' Naomi chỉ tay về phía Suzuki, ``\vie{đã gọi anh là `Onii-sama' rồi. Em không muốn gọi giống chị ấy. Còn chị đằng kia,}'' em liếc sang Kaigou, ``\vie{thì gọi anh là Kuon-kun. Em tôn trọng anh là người lớn, nên em đã chọn thêm hậu tố `-chan' để phân biệt.}''

Thực ra, ở thời kỳ cổ đại của Nhật Bản, hậu tố ``-chan'' chưa hề tồn tại, nó chỉ là một biến thể ngọng nghịu của ``-san'' dành cho trẻ con được dùng sau này. Nhưng có vẻ như gói ngôn ngữ mà Game tải vào đã cung cấp cho cô bé những lựa chọn ngữ pháp hiện đại nhất, và em đã quyết định sử dụng cách xưng hô thân mật nhất (chỉ kém việc không dùng hậu tố nào) để tạo ra một sự gắn kết đặc biệt với tôi.

``\vie{Thế còn hai chị ấy thì em gọi như thế nào?}'' tôi tò mò.

``\vie{Suzuki-onee-chan và Kaigou-onee-chan ạ.}''

Nghe xong, Suzuki vỡ òa, lao tới ôm chầm lấy cô bé. ``\vie{Ôi trời ơi, dễ thương quá đi mất! Em có thể gọi chị như thế hàng ngày luôn được không!?}''

Kaigou cũng đỏ mặt, gãi gãi má nhưng miệng thì mỉm cười không giấu giếm: ``\vie{Hừm\dots\ cũng không tệ. Nghe lọt tai đấy.}''

Sau một hồi loay hoay thoát khỏi vòng tay siết chặt của Suzuki, Naomi đưa ánh mắt nhìn ra ngoài cửa sổ, nơi ánh nắng chan hòa đang phủ kín những mái nhà và con phố lát đá của Aogami.

``\vie{Thế giới ngoài kia\dots\ bây giờ trông như thế nào vậy?}'' cô bé lẩm bẩm, một sự khao khát không lời hiện rõ trong đáy mắt lục bảo. ``\vie{Nó có khác nhiều so với thời của em không?}''

Câu hỏi ngây ngô ấy khiến tôi chợt nhận ra. Kể từ khi bắt đầu ``đóng đô'' tại Aogami để thực hiện nhiệm vụ giải nguyền, chúng tôi đã luôn ở trong trạng thái cảnh giác, làm việc bán thời gian, điều tra oán linh và cắm mặt vào những bí ẩn. Chúng tôi chưa bao giờ thực sự có một ngày nghỉ đúng nghĩa, một buổi đi chơi thư giãn cùng nhau.

Tôi thở dài, mỉm cười rồi đứng dậy khỏi ghế. ``\vie{Em hỏi không thôi thì khó hình dung lắm. Hôm nay đẹp trời, mọi người có muốn ra ngoài một chuyến không?}''

\pov{Hikawa Suzuki}

Tôi hào hứng chuẩn bị đồ đạc ngay khi Onii-sama đề nghị một buổi đi chơi. Lâu lắm rồi cả ba chúng tôi mới có một ngày thảnh thơi thực sự.

Nhưng trước tiên, chúng tôi có một vấn đề cần giải quyết: trang phục của Naomi-chan.

Bộ kimono mộc mạc mà em đang mặc tuy rất đẹp và hợp với nét tĩnh lặng của em, nhưng nó hoàn toàn không phù hợp cho một buổi đi dạo trong thế giới hiện đại. Hơn nữa, thời tiết đầu tháng Sáu đã bắt đầu mang theo hơi nóng của mùa hè, mặc một bộ đồ nhiều lớp như vậy sẽ rất bức bối. Tôi quyết định là sau khi mua đồ mới, tôi sẽ giặt giũ và cất cẩn thận bộ kimono đó, coi như một kỷ vật thiêng liêng về nguồn cội của em.

Sau khi thay đồ xong, cả bọn chúng tôi xuống phố.

Nắng cuối tuần đầu hè chưa quá gắt, chỉ đủ để lấp lánh trên mái ngói đã ngả màu thời gian. Phố Aogami vào sáng thứ Bảy như một bản giao hưởng nhỏ: tiếng chuông xe đạp leng keng, mùi bánh taiyaki nóng thoang thoảng từ quán góc phố, và những tia nắng rọi xiên qua tán lá tươi xanh mướt. Tôi hít một hơi thật sâu, mùa hè ở đây khác hẳn ở thế giới cũ của tôi, không có khói bụi, chỉ toàn gió núi và hương cỏ non. Thị trấn Aogami hiện ra dưới tầm mắt: những mái nhà ngói xếp lớp, xa xa là cánh đồng trà xanh bậc thang, và trên cao, đường tàu một ray vắt ngang sườn núi. Gió thổi từ thung lũng mang theo hương trà non và mùi đất nung ấm. Tôi chợt nhận ra, trong mắt Naomi, thị trấn này không chỉ là nơi ở mới, đó là cả một thế giới đang chờ em khám phá, từng milimet, từng hơi thở.

Điểm đến đầu tiên của bốn người chúng tôi là khu phố mua sắm dưới chân núi. Mục tiêu của chúng tôi lúc này là sắm cho cô bé những bộ trang phục mới để có thể sống trong thế giới này.

``Chị thấy cái này thế nào?'' tôi giơ lên một chiếc váy mùa hè màu trắng tinh khôi điểm xuyết họa tiết hoa bỉ ngạn nhạt màu. ``Chất liệu mỏmg, rất thoáng mát cho mùa hè. Đi kèm với mũ cói trùm nắng và đôi sandal quai hậu này là chuẩn bài!''

Chị tôi đứng bên cạnh, cầm theo vài bộ quần áo mặc nhà gọn nhẹ gồm áo phông và quần short thun, gật gù đồng ý: ``Duyệt. Ở nhà cứ mặc thế này cho thoải mái. Nhìn con bé mặc kimono làm chị cứ tưởng mình đang xem phim truyền hình Taiga lịch sử ấy.''

Naomi đứng trong phòng thử đồ bước ra. Chiếc váy tôn lên vóc dáng nhỏ nhắn và làm nổi bật nước da trắng của em. Kết hợp với mái tóc đen tết bím và đôi mắt xanh lục bảo, trông em hệt như một cô búp bê sứ vừa bước ra từ một bức tranh mùa hè rực rỡ. Dù có vẻ hơi bối rối với việc để hở cánh tay và đôi chân (điều tối kỵ ở thời phong kiến của em), nhưng ánh mắt em lại lấp lánh sự thích thú.

``Naomi-chan, trông em xinh lắm!'' tôi vỗ tay, mắt long lanh. ``Đúng kiểu thiên thần mùa hè!''

``Ừ, váy trắng hợp với da và tóc đen tuyền kia,'' Onee-sama gật gù, khẽ xoay em một vòng. ``Không trang điểm mà vẫn bật ra nét thơ.''

Chúng tôi liếc về phía Onii-sama, chờ đợi một câu khen ngợi từ anh. Nhưng đáp lại, anh chỉ đứng đó, hai tay đút túi, nhìn thoáng qua rồi khẽ gật gù, im bặt - kiểu không khen mà cũng không chê.

Naomi bối rối, mắt lục bảo chớp chậm, nhỏ giọng: ``\dots Có gì không ổn sao ạ?''

Tôi huých khuỷu tay anh tôi, thì thầm: ``Onii-sama, anh nói gì đi chứ!''

Onee-sama cũng cà khịa: ``Kuon-kun, im lặng kiểu đó cô bé tưởng mình mặc xấu đấy.''

Kuon khẽ nhíu mày, lẩm nhẩm: ``\dots Trông em cũng ổn ổn đấy. ''

Cô bé nghe vậy, tai đỏ ửng, cúi xuống vuốt vạt váy; nhưng đôi mắt vẫn hướng về Kuon, chờ thêm lời. Thấy anh không nói thêm, tôi vội cười trừ: ``Ý anh ấy là rất được đó, em đừng lo!''

Onee-sama thở dài, xoa đầu em: ``Thôi, coi như tên kia gật đầu chấp thuận rồi nhé. Đi thôi, còn nhiều đồ phía trước!''

Cứ như thế, chúng tôi tiếp tục thử đồ cho Naomi, từ những bộ trang phục ngủ đơn giản mà ấm áp, cho tới những bộ đồ hiện đại hơn (mặc dù hơi lạc quẻ so với tính cách trầm ngâm của em ấy). Mỗi lần thử đồ, em ấy đều hiện lên vẻ ngây thơ và dễ thương theo nhiều cách khác nhau.

``Tuyệt vời! Lấy hết chỗ này đi!'' tôi reo lên.

``Này, này, Suzuki, em lấy hết thì nhà mình cất ở đâu?'' Onii-sama hỏi tôi với giọng nửa châm chọc, nửa ái ngại.

``Chị biết là em rất hứng thú với em ấy,'' Onee-sama tiếp lời, ``nhưng em cũng phải kiềm chế lại bản thân đi chứ. Mua nhiều quá mà không mặc lại phí tiền đấy.'' Tôi chỉ cười trừ với phản ứng của hai anh chị.

Đứng ở quầy thu ngân, anh tôi thản nhiên rút ví thanh toán toàn bộ số quần áo mới. Kể ra thì, số tiền anh ấy kiếm được từ công việc bán thời gian ở cửa hàng tiện lợi những ngày qua khá rủng rỉnh. Anh ấy từng nói rằng số tiền này dù sao cũng không mang về thế giới của anh được, nên tiêu hết ở đây cũng chẳng sao. Không ngờ, nó lại trở thành ``quỹ nuôi trẻ'' bất đắc dĩ.

Rời khỏi cửa hàng, chúng tôi đi dạo dọc theo những con phố trung tâm Aogami, với bộ đồ đầu tiên em ấy ướm thử, trong khi bộ kimono mà em mặc trước đó được tôi cất cẩn thận trong đó.

Thế giới trong con mắt của một cô bé đến từ 1300 năm trước thực sự là một kỳ quan. Mỗi khi một chiếc xe ô tô chạy ngang qua, cô bé lại dừng bước, đôi mắt diễu theo chuyển động của cỗ máy kim loại với vẻ ngạc nhiên không giấu giếm. Khi nghe tiếng chuông reng reng của chiếc xe điện chạy dưới lòng đường, em hơi co vai lại, núp nép sau lưng tôi, nhưng rồi lại thò đầu ra nhìn với vẻ say mê.

``Thế giới bây giờ\dots\ không còn dùng xe ngựa kéo nữa sao?'' Em hỏi nhỏ, ngước nhìn lên cột đèn giao thông đang chuyển từ đỏ sang xanh.

``Không đâu em,'' tôi mỉm cười giải thích, nắm chặt bàn tay nhỏ xíu của em. ``Bây giờ mọi thứ chạy bằng điện và động cơ. Dù đường phố ồn ào hơn một chút, nhưng em thấy đấy, nó không hề đáng sợ đúng không?''

``Vâng\dots'' Em khẽ gật đầu, ngắm nhìn những tòa nhà bê tông và cửa kính phản chiếu ánh mặt trời. ``Rất nhộn nhịp. Cảm giác như\dots\ mọi người đều đang sống hết mình vậy.''

Chúng tôi tiếp tục dạo quanh khu phố thị sầm uất.

Khi đi ngang qua một cửa hiệu thuốc lớn, Naomi bỗng khựng lại. Em ngước nhìn bảng quảng cáo điện tử treo trước cửa, màn hình LED nhấp nháy chuyển từ màu tím sang hồng rồi xanh lục như một khối pha lê khổng lồ biết tự phát sáng. Đôi mắt lục bảo của em mở to, phản chiếu dòng chữ ``Mở cửa 24h'' đang chạy ngang màn hình.

``Ánh sáng\dots\ đầy sắc màu kìa? Không cần đốt nến mà vẫn sáng sao?'' em thì thầm, tay vô thức nắm chặt lấy vạt váy trắng tinh.

Onee-sama cười khúc khích, dịu dàng xoa đầu em: ``Đó là đèn LED đấy em ạ, chạy bằng điện năng chứ không phải ma thuật hay nến đâu. Đi thôi, thế giới này còn nhiều điều hay ho hơn thế nhiều.''

Tiếp tục băng qua con dốc đá cũ kỹ lát gạch đỏ, chúng tôi bắt gặp mấy bà cụ đang trải bạt bán những rổ dâu tây núi chín mọng. Bất chợt, một cậu bé lao vọt qua chúng tôi trên chiếc xe trượt scooter nhỏ, phần tay lái phát ra tiếng nhạc điện tử chói tai nhấp nháy đèn. Naomi giật mình, vội núp lùi về sau lưng tôi, nhưng rồi lập tức thò đầu ra, đôi mắt long lanh dõi theo cậu nhóc: ``Chiếc xe kêu như chim hót!''

Tôi không nhịn được cười, vỗ vỗ vai em: ``Đúng rồi, chim điện đấy. Ở thời đại này người ta nuôi nhiều `chim' bằng pin lắm.''

Đi dạo thêm một lúc dưới cái nắng đầu hè, cái nóng oi ả bắt đầu làm cả bọn thấm mệt. Onii-sama, với tư cách là ``nhà tài trợ'' hào phóng nhất ngày hôm nay, chỉ tay về phía một quán cà phê tráng miệng ở góc phố: ``Nghỉ chân chút nhé. Trông ba người đều có vẻ khát rồi.''

Bên trong quán được lắp điều hòa mát rượi khiến Naomi khẽ rùng mình khoan khoái. Khi đi ngang qua quầy pha chế, máy đánh sữa espresso rít lên từng hồi, tỏa ra những làn hơi nước trắng xóa. Em úp hai bàn tay nhỏ xíu lên tấm kính trong suốt, mũi ưỡn ra áp sát vào lớp kính nhìn lớp bọt sữa đang phồng lên.

``Sương mù\dots\ bị nhốt trong cốc kìa?'' em ngây ngô chỉ trỏ.

Onii-sama gõ nhẹ lên trán em, kéo em về bàn: ``Không phải sương đâu, hơi nước đấy. Chờ chút, anh sẽ gọi cho em một thứ còn kỳ diệu hơn.''

Một lát sau, nhân viên phục vụ bưng ra bốn ly kem parfait trà xanh khổng lồ, bên trên phủ đầy đậu đỏ, thạch mochi, và một lớp kem vani trắng muốt. Naomi trợn tròn mắt nhìn ly kem cao ngất ngưởng, cẩn thận dùng chiếc thìa dài múc một miếng kem nhỏ đưa lên miệng. Khoảnh khắc hơi lạnh chạm vào đầu lưỡi, em khẽ rùng mình, hai mắt mở to như vừa phát hiện ra châu lục mới.

``Lạnh\dots\ lạnh quá! Nhưng mà ngọt! Tuyết giữa mùa hè sao?'' em thốt lên, hai má ửng hồng vì thích thú. Tốc độ xúc kem của em tăng dần, hoàn toàn rũ bỏ vẻ điềm tĩnh ban đầu để trở về đúng với bản tính của một cô bé tám tuổi. ``Ngon quá đi mất\dots''

Ăn xong, chúng tôi lại rảo bước dọc theo con kênh nhỏ dẫn về phía công viên. Chiếc cầu đá vòm cong phủ đầy rêu phong có lẽ là thứ gợi cho em nhiều sự quen thuộc nhất. Đứng dưới chân cầu nhìn xuống dòng nước trong vắt, mấy con cá Koi màu đỏ tươi pha trắng đang lượn lờ bơi lội. Naomi ngồi xổm xuống, đưa ngón tay trỏ khẽ khuấy mặt nước. Đám cá Koi tưởng có đồ ăn liền há miệng quẫy đuôi xúm lại.

Em reo lên, nụ cười hiếm hoi lại nở rộ: ``Là ngựa biển có cánh! Đẹp quá đi mất!''

Onee-sama cười nắc nẻ, ngồi xổm xuống cạnh em: ``Đó là cá Koi, không phải ngựa. Nhưng mà\dots\ nếu em muốn hiểu cái vây là `cánh' thì cũng được.''

Sau khi đi dạo quanh khu phố thị, chị tôi đề xuất chúng tôi tới công viên trung tâm. Nơi đây từng được Inari - khi anh chị vẫn còn là ``Kyuen'' và ``Saikyou'' - nhắc tới như một biểu tượng của những ngày tháng yên bình chưa bị bóng tối bao phủ. Và quả thật, công viên rất rộng rãi với những thảm cỏ xanh mướt, cây xanh rợp bóng và khu vui chơi với đầy đủ cầu trượt, xích đu và hố cát.

Chúng tôi quyết định để Naomi trải nghiệm mọi thứ.

Bắt đầu bằng trò cầu trượt. Ban đầu, em ấy chỉ e dè đứng nhìn từ dưới lên, không dám bước lên những bậc thang nhựa. Tôi liền nắm lấy tay em, kéo em leo lên đỉnh cầu trượt. Mặc kệ việc mình đã mười lăm tuổi, lớn đùng rồi chứ chẳng nhỏ nhắn gì, tôi ôm trọn Naomi vào lòng rồi cùng trượt xuống.

``Kyaaa!'' Tiếng hét pha lẫn tiếng cười giòn tan của em vang vọng khi cả hai chúng tôi hạ cánh xuống thảm cát mềm.

Kế đến là xích đu. Naomi ngồi lọt thỏm trên một chiếc, còn tôi chiếm luôn chiếc bên cạnh.

``Onii-sama, Onee-sama! Đẩy đi nào!'' tôi quay đầu lại hối thúc.

Onii-sama thở hắt ra, lầm bầm: ``Cái con bé này, rốt cuộc là đang dẫn trẻ con đi chơi hay tự mình đi chơi vậy?'' Dẫu phàn nàn là thế, tay anh vẫn nhịp nhàng đẩy chiếc xích đu của Naomi lên, trong khi Onee-sama thì đẩy bổng tôi lên cao tít.

Naomi nhắm nghiền mắt lại, cảm nhận những cơn gió tạt vào mặt mỗi khi xích đu lao tới trước. ``Giống như\dots\ đang bay vậy!'' em reo lên, hai bím tóc tung bay phấp phới.

Không dừng lại ở đó, khu công viên rộng lớn với thảm cỏ xanh lập tức trở thành sân chơi cho trò đuổi bắt. Tôi làm ``quỷ'', đuổi theo Naomi đang lạch bạch chạy phía trước bằng đôi sandal mới. Onee-sama chạy theo yểm trợ chặn đường em ấy, trong khi Onii-sama thì chỉ đứng khoanh tay dựa lưng vào gốc cây, cười khùng khục nhìn ba chị em gái chạy tán loạn. Cảnh tượng ồn ào ấy thu hút ánh nhìn của vài người đi dạo trong công viên, nhưng chúng tôi chẳng ai bận tâm.

Nụ cười mỉm của em xuất hiện nhiều hơn, không còn nét u uất của một linh hồn bị giam cầm trong quá khứ. Nhìn em vui vẻ rạng rỡ, trong lòng tôi trào dâng một cảm giác ấm áp khó tả. Suzuki này\dots\ cuối cùng cũng có thể bảo vệ được một người rồi.

\ct

Chúng tôi tiếp tục dạo bước qua những ngóc ngách của thị trấn—tôi bất giác thấy lòng mình nhẹ nhõm đến lạ. Bao tháng ngày dài đằng đẵng chỉ biết cắm mặt vào hồ sơ oán linh, rồi thức trắng đêm lo cho từng vết nứt của lời nguyền, nay được thay bằng tiếng cười giòn tan của Naomi mỗi khi cô bé giật mình hay reo lên khi thấy một thứ gì đó mới lạ mà ở thời phong kiến chưa được thấy. Những phản ứng ấy như phản chiếu chính bản thân tôi ngày đầu bước vào Aogami, một thế giới mới mẻ, rực rỡ và đầy sức sống nhưng có gì đó hơi giả tạo xen lẫn trách nhiệm, còn em, đứa trẻ vừa tái sinh từ 1300 năm bóng tối, chỉ đơn thuần ngỡ ngàng rồi mỉm cười. Tôi nhận ra mình đang sống lại ký ức tuổi thơ qua đôi mắt lục bảo ấy: mỗi chiếc xe điện reng reng qua đường, mỗi tòa nhà kính phản chiếu hoàng hôn, đều trở thành phép màu nho nhỏ.

Có lẽ đây mới đúng là ``du hí,'' không chỉ đi để khám phá, mà để nhắc nhau rằng sau bao mùa mưa gió, ta vẫn còn quyền được vui, được thấy cuộc đời rạng rỡ như ánh nắng đầu hè đang rải xuống phố Aogami.

Cuối buổi chiều, khi hoàng hôn nhuộm màu cam rực rỡ lên bầu trời Aogami, bốn chúng tôi ngồi nghỉ trên một băng ghế dài nhìn ra hồ nước giữa công viên. Gió thổi mơn man, mang theo hơi mát của nước. Naomi ngồi kẹp giữa tôi và Onee-sama, trong khi Onii-sama đứng dựa lưng vào cột đèn cạnh đó, nhâm nhi một lon cà phê ướp lạnh.

``\vie{Thế nào Naomi, cảm nhận ngày đầu tiên trở thành người hiện đại?}'' anh tôi lên tiếng hỏi, phá vỡ bầu không khí tĩnh lặng.

Naomi đưa tay chỉnh lại vành mũ cói, ánh mắt xa xăm nhìn về phía mặt trời lặn.

``\vie{Nó\dots\ rất ấm áp,}'' giọng em nhẹ nhàng, nhưng lại mang một độ sâu kỳ lạ. ``\vie{Không còn sương mù hay mùi đất. Không còn oán hận hay đau khổ. Em có thể cảm nhận được mùi của cỏ, tiếng chim hót và cả hơi ấm từ bàn tay của mọi người. Ở thời của em, hòa bình là một thứ gì đó vô cùng xa xỉ. Còn ở đây\dots\ mọi thứ đều sống động và rạng rỡ.}''

Em quay sang nhìn ba chúng tôi, đôi mắt xanh lục lấp lánh phản chiếu ráng chiều. ``\vie{Cảm ơn các anh chị. Vì đã đưa em tới thế giới này.}''

Nghe những lời chân thành từ một cô bé vừa trải qua cả ngàn năm chịu đựng sự ngột ngạt của lời nguyền, trái tim tôi như mềm ra.

Onee-sama tựa cằm lên tay, bật cười một tiếng rồi buông một câu đùa bâng quơ: ``\vie{Trời ạ, nhìn chúng ta bây giờ đi. Một anh chàng lớn tuổi trả tiền cho mọi thứ, hai cô nàng xách đồ dạo phố, và một đứa trẻ lăng xăng khám phá thế giới. Nhìn từ xa, bốn người đi chơi chung với nhau chẳng khác gì một gia đình nhỏ đang đi dã ngoại cuối tuần ghê.}''

Câu đùa của Onee-sama lập tức tạo ra những phản ứng dây chuyền thú vị.

Onii-sama đang uống dở ngụm cà phê thì khựng lại, ho sặc sụa rồi nhíu mày lườm chị tôi như thể cô vừa nói điều gì đó xúc phạm lắm. Còn tôi thì cảm nhận rõ hai bên má mình đang nóng ran lên. Bị ghép vào hình ảnh ``gia đình'' một cách cụ thể như thế này\dots\ khiến tôi không tránh khỏi một cảm giác gì đó bối rối.

Nhưng điều ngạc nhiên nhất là Naomi. Cô bé không hề tỏ ra ngượng ngùng hay khó hiểu. Ngược lại, em ấy khẽ mỉm cười - một nụ cười hiền hậu và rạng rỡ nhất từ sáng tới giờ.

``\vie{Một gia đình sao\dots}'' em thầm thì, ngả đầu tựa vào vai tôi. ``\vie{Em cũng mong có một gia đình ấm áp như thế này.}''

Câu chốt đầy cảm xúc ấy đánh gục hoàn toàn chút sự phòng bị cuối cùng của ba chúng tôi. Tôi cảm thấy trái tim mình như bị bóp nghẹt. Một đứa trẻ đã chờ đợi hơn một thiên niên kỷ chỉ để được nói ra hai tiếng ``gia đình'' một cách giản dị như thế. Trong mắt em, mái ấm không phải là ngôi nhà lớn, không phải là của cải vật chất, mà chỉ đơn thuần là bốn con người ngồi cạnh nhau dưới hoàng hôn, chia sẻ tiếng cười và cảm giác được yêu thương. Tôi bất giác siết nhẹ bàn tay đang nắm lấy em - như thể chỉ cần buông lỏng một giây, em sẽ lại rơi về cô đơn ngàn năm kia.

Gió vẫn thổi nhẹ qua công viên, cuốn đi những cánh hoa bay lác đác. Trong khoảnh khắc ấy, tôi biết chắc chắn một điều: thị trấn Aogami đã thực sự tìm lại được sự bình yên, và chúng tôi cũng đã tìm được một mảnh ghép mới cho mái ấm của mình. Trong ánh hoàng hôn rực rỡ, với Naomi nép vào lòng, tôi bỗng nhận ra: chúng tôi không chỉ đơn thuần là những người bảo vệ hay giải cứu. Hơn cả một sứ mệnh, hành trình này đã trở thành một điều gì đó ấm áp và sâu sắc hơn. Chúng tôi, những mảnh ghép lạc lõng tìm thấy nhau giữa dòng chảy thời gian, đang cùng nhau tạo nên một mái ấm mới: một gia đình nhỏ nơi sự cô đơn tan biến và yêu thương đâm chồi nảy lộc. Chúng tôi không chỉ đang bảo vệ một linh hồn bé nhỏ khỏi lời nguyền, chúng tôi đang dựng lại cho em điều mà thời gian và đau thương đã cướp mất - quyền được tin rằng mình xứng đáng có một nơi để về.

Có lẽ, đây mới chính là phép màu thực sự mà Aogami cần, và cũng là điều trái tim tôi đã tìm kiếm bấy lâu.

\end{document}